\documentclass[10pt,a4paper]{amsart}
\usepackage[margin=19mm]{geometry}
\usepackage{amsmath,amssymb,mathtools}
\usepackage[scr=rsfs,cal=euler]{mathalfa}
\usepackage{xfrac}
\usepackage[unicode,hidelinks,bookmarks=false]{hyperref}
\newcommand{\ie}{i.e.\ }
\newcommand{\cf}{cf.\ }
\newcommand{\resp}{respectively\ }
\newcommand{\Subsection}{\subsection}
\providecommand{\nobreakdash}{\nobreak\hbox{-}}
\setlength{\emergencystretch}{2em}
\multlinegap0pt
\usepackage{amsthm}
\usepackage{mathtools}
\usepackage{amsmath,amssymb,amsfonts}
\usepackage{graphicx}
\usepackage{bm}
\usepackage{enumitem}
\usepackage{mathrsfs}
\usepackage{xcolor}
\usepackage{tikz}
\usetikzlibrary{arrows.meta,calc,positioning,patterns,decorations.pathmorphing,decorations.markings}
\usepackage{tikz-cd}
\providecommand{\CC}{\mathbf{C}}
\DeclareMathOperator{\rank}{rank}
\numberwithin{equation}{section}
\theoremstyle{plain}
\newtheorem{theorem}{Theorem}[section]
\newtheorem{lemma}[theorem]{Lemma}
\newtheorem{proposition}[theorem]{Proposition}
\newtheorem{corollary}[theorem]{Corollary}
\newtheorem{thmAlph}{Theorem}
\theoremstyle{definition}
\newtheorem{remark}[theorem]{Remark}
\newtheorem{example}[theorem]{Example}
\theoremstyle{plain}
\newtheorem{prop}[theorem]{Proposition}
\newtheorem{teo}[theorem]{Theorem}
\newtheorem{cor}[theorem]{Corollary}
\theoremstyle{definition}
\newtheorem{defin}[theorem]{Definition}
\title[Ramification of the moduli map for hyperplane sections]{Ramification of the Moduli Map for Hyperplane Sections of K3s and Hypersurfaces: the incidence bound for the restricted differential (Lemma 3.3)}
\author[D. Faro, F. Gounelas]{Dario Faro, Frank Gounelas}
\date{Source: arXiv:2609.29501v1 (CC BY 4.0)}
\begin{document}
\maketitle
\vspace{1em}
\noindent\textit{Source and licence.} Ramification of the moduli map for hyperplane sections of K3s and hypersurfaces. Version arXiv:2609.29501v1 (CC BY 4.0), distributed under the Creative Commons Attribution 4.0 International licence, http://creativecommons.org/licenses/by/4.0/, as recorded in the arXiv OAI licence metadata for this exact version and shown on the abstract page https://arxiv.org/abs/2609.29501v1. This item is an editorial re-presentation of the paper's incidence bound for the linear map associated with the restricted differential of a first-order deformation, printed as Lemma 3.3. A full rights notice is delivered at licenses/arxiv-2609.29501-CC-BY-4.0.txt.
\noindent\textit{Editorial scope.} Counted unit: the paper's incidence bound for the linear map associated with the restricted differential of a first-order deformation, printed as Lemma 3.3 (source0.tex lines the statement at lines 665--700 and the proof at lines 702--922), delivered together with the paper's complete original proof of it as the single primary proof body. No mathematics is written, repaired or re-derived: statement and proof are the paper's own text line for line, in the paper's own notation. Every cross-reference inside the retained spans resolves inside the delivered document. Printed numbers are the paper's own: the wrapper restores the paper's counters before the retained material, so the counted statement prints as Lemma 3.3, the paper's own number for it, and the numbered displays of the retained spans carry the paper's own section-relative equation numbers. Omitted from this package and not redistributed: the paper's preamble and front matter as such (of which only the title and the author names are reproduced in the wrapper); the sections, lemmas, theorems and examples that lie outside the retained spans; and the paper's bibliography with its bib data file. No citation command occurs in any retained span, so the delivered document carries no bibliography and no reference or citation command of the delivered text was rewritten or adapted. Printed slips kept literal, among them the paper's own wording and its own choice of symbols. Layout and class: the delivered wrapper is the lane's pinned amsart editorial wrapper at 10pt on A4, to which this package adds the paper's own theorem-like declarations and the shorthand macros its retained text uses, all copied verbatim from the paper's source. No publisher class file, no publisher style file and no upstream script is redistributed. Assets: no retained span of this package contains a figure environment or a graphics-inclusion command, no image or artwork file exists in or is redistributed with this package, and the manifest and the delivery evidence both carry an empty asset-dependency list. A full rights notice is delivered at licenses/arxiv-2609.29501-CC-BY-4.0.txt.
\setcounter{section}{3}
\setcounter{theorem}{2}
\begin{lemma}\label{lem:incidence-bound}
Let \(n\ge3\) and \(d\ge4\). For
\[
D=\sum_{i=0}^{n-1}A_i\partial_i,
\qquad A_i\in V_{n,1},
\]
consider the linear map
\[
\begin{tikzcd}[row sep=small]
V_{n,d}\oplus V_{n,d-1}
  \arrow[r,"M_D"]
& V_{n,d},\\
(g,h)\arrow[r,mapsto]
&D(g)-x_{n-1}h,
\end{tikzcd}
\]
and the  incidence variety
\[
\mathcal I_{n,d}:=
\left\{
(g,h,D)\in
V_{n,d}\times V_{n,d-1}\times(V_{n,1})^n
\ \middle|\ M_D(g,h)=0
\right\}.
\]
Then
\[
\dim\mathcal I_{n,d}\le \dim V_{n,d}+n^2-n.
\tag{1}\label{eq:I-general-bound}
\]
In the exceptional case \((n,d)=(3,4)\), one has the stronger estimate
\[
\dim\mathcal I_{3,4}\le19.
\tag{2}\label{eq:I-exceptional-bound}
\]
\end{lemma}

\begin{proof}
We will estimate the dimension of \(\mathcal I_{n,d}\) using the projection
\[
\begin{aligned}
\pi\colon\mathcal I_{n,d}&\longrightarrow(V_{n,1})^n,\\
(g,h,D=\textstyle\sum A_i\partial_i)&\longmapsto(A_0,\dots,A_{n-1}).
\end{aligned}
\]
The fibre of \(\pi\) over \(D\) is exactly \(\ker M_D\). We first describe in a convenient way the kernel of $M_D$. The restriction of \(M_D\) to \(0\oplus V_{n,d-1}\) maps isomorphically
onto
\[
x_{n-1}V_{n,d-1}\subseteq V_{n,d}.
\]
Thus this part contributes \(\dim V_{n,d-1}\) to the rank. The remaining
contribution is the dimension of the image in $V_{n,d}/x_{n-1}V_{n,d-1} \simeq V_{n-1,d}$, that is we need to compute the dimension of the image of the linear map
\[
\begin{tikzcd}[row sep=small]
\phi_D\colon V_{n,d}\oplus V_{n,d-1}
  \arrow[r]
&V_{n,d}/x_{n-1}V_{n,d-1}\simeq V_{n-1,d},\\
(g,h)\arrow[r,mapsto]
&D(g)|_H,
\end{tikzcd}
\]
where $H:=\{x_{n-1}=0\}$.
 Set
\[
\delta_D:=
\sum_{i=0}^{n-2}(A_i|_H)\partial_i
\in \operatorname{Der_{\CC}{}}(V_{n-1})_0,
\qquad
a_D:=A_{n-1}|_H\in V_{n-1,1}.
\]
Here $\operatorname{Der}_{\CC}(V_{n-1})_0$ denotes the space of $\CC$-linear derivations of $V_{n-1}$ that preserve degree. Restriction to $V_{n-1,1}$ gives an isomorphism $\operatorname{Der}_{\CC}(V_{n-1})_0\simeq\operatorname{End}(V_{n-1,1})$, whose inverse is
\begin{equation}
\label{iso derivation}
\operatorname{End}_{\CC}(V_{n-1,1})
\xrightarrow{\sim}
\operatorname{Der}_{\CC}(V_{n-1})_0,
\qquad
T\longmapsto
D_T:=\sum_{i=0}^{n-2}T(x_i)\partial_i.
\end{equation}
In the following we switch freely between these two descriptions. Writing
\[
g=b_0+x_{n-1}b_1+x_{n-1}^2b_2+\cdots,
\qquad
b_k\in V_{n-1,d-k},
\]
we obtain
\[
D(g)|_H=\delta_D(b_0)+a_Db_1.
\]
Consider the linear map
\[
\Phi_{\delta_D,a_D}\colon V_{n-1,d} \oplus V_{n-1,d-1}  \rightarrow V_{n-1,d}, \quad  \quad (b_0,b_1) \rightarrow \delta_D(b_0) 
 +a_D(b_1)
\]
and set $r(\delta_D,a_D):=\rank\Phi_{\delta_D,a_D}$. To be more precise here, we are thinking $\delta_D \in \operatorname{Der}_{\CC}(V_{n-1})_0$ (via \eqref{iso derivation}), and restricting its action on $V_{n-1,d}$. Since the map
\[
V_{n,d}\longrightarrow V_{n-1,d}\oplus V_{n-1,d-1},
\qquad
g\longmapsto(b_0,b_1),
\]
is surjective, the rank of $\phi_D$ is \(r(\delta_D,a_D)\). Then
\[
\rank M_D=\dim V_{n,d-1}+r(\delta_D,a_D)
\tag{3}\label{eq:rank-MD-two-parts}
\]
and
\[
\begin{aligned}
\dim\ker M_D
&=\dim V_{n,d}+\dim V_{n,d-1}-\rank M_D\\
&=\dim V_{n,d}-r(\delta_D,a_D).
\end{aligned}
\tag{4}\label{eq:kernel-MD-two-parts}
\]
We write now the parameter space $(V_{n,1})^n$ as a disjoint union of locally closed subvarieties for which we are able to   estimate \(r(\delta_D,a_D)\). For this purpose, consider the linear map
\begin{equation}
\label{parametersmap}
\begin{aligned}
(V_{n,1})^n
&\longrightarrow \operatorname{End}(V_{n-1,1})\oplus V_{n-1,1}\\
D&\longmapsto(\delta_D,a_D).
\end{aligned}
\end{equation}
It is surjective and has \(n\)-dimensional kernel. We decompose the parameter space $(V_{n,1})^n$ into the inverse images of the three loci in $\operatorname{End}(V_{n-1,1})\oplus V_{n-1,1}$ defined by $a\ne0$, by $a=0$ and $\delta\ne0$, and by $a=\delta=0$. If $a\ne0$, multiplication by $a$ identifies $V_{n-1,d-1}$ with the subspace $aV_{n-1,d-1}\subset V_{n-1,d}$, so $r(\delta,a)\ge \dim V_{n-1,d-1}$; using also \eqref{parametersmap} we find that  the inverse image of this locus in $(V_{n,1})^n$ has dimension $n(n-1)+n=n^2$. If $a=0$ and $\delta\ne0$, then $r(\delta,0)\ge1$: indeed, choose $\ell\in V_{n-1,1}$ with $\delta(\ell)\ne0$, and then $\delta(\ell^d)=d\ell^{d-1}\delta(\ell)\ne0$. The inverse image of this  locus has dimension $n+(n-1)^2=n^2-n+1$. Finally, if $a=\delta=0$, then $r(\delta,a)=0$ and the corresponding locus has dimension $n$. 

Consequently, the dimensions of the parts of $\mathcal I_{n,d}$ lying  over these three strata are bounded from above respectively by
\[
   \dim V_{n,d}+n^2-\dim V_{n-1,d-1},\qquad
   \dim V_{n,d}+n^2-n,\qquad
   \dim V_{n,d}+n.
\]
Since $\dim V_{n-1,d-1}\ge n$, all three dimensions are at most $\dim V_{n,d}+n^2-n$. Hence
\[
\dim\mathcal I_{n,d}\le\dim V_{n,d}+n^2-n.
\]
This proves \eqref{eq:I-general-bound}. Now we consider the case $(n,d)=(3,4)$. In this situation
\[
\dim V_{3,4}=15,
\qquad
\dim V_{2,3}=4,
\qquad
\dim V_{2,4}=5.
\]
We want to give a better bound for $\operatorname{dim}(\mathcal{I}_{3,4})$. For $(\delta,a) \in \operatorname{End}(V_{2,1}) \oplus V_{2,1}$ we consider as before
\[
\Phi_{\delta,a}\colon V_{2,4} \oplus V_{2,3} \rightarrow V_{2,4}.
\]


As above, if $a\ne0$, then $aV_{2,3}\subset\operatorname{im}\Phi_{\delta,a}$ has dimension $4$. Since $\dim V_{2,4}=5$, the rank of $\Phi_{\delta,a}$ is therefore either $4$ or $5$. Thus we further stratify the locus $\{a \neq 0\}\subset \operatorname{End}(V_{2,1}) \oplus V_{2,1} $ as $T_4 \sqcup T_5$,
where for $r\in\{4,5\}$,  we set
\[
T_r:=
\left\{
(\delta,a)\in
\operatorname{End}(V_{2,1})\times
\bigl(V_{2,1}\setminus\{0\}\bigr)
\ \middle|\
\rank\Phi_{\delta,a}=r
\right\}.
\]
We denote by $S_4$ and $S_5$ the inverse images via \eqref{parametersmap}, namely
\[
   S_r:=\{D\in(V_{3,1})^3\mid a_D\ne0,\
   \rank\Phi_{\delta_D,a_D}=r\}.
\]
Since $\dim S_5\leq \operatorname{dim}(V_{3,1})^3=9$, we have
\[
   \dim\{(g,h,D) \in \mathcal{I}_{3,4}\mid D\in S_5\}
   \le\dim V_{3,4}+9-5=\dim V_{3,4}+4=19.
\]
Now we need to bound the dimension of $S_4$. On $S_4$ the composite
\[
\begin{tikzcd}
V_{2,4} \arrow[r, "\delta"] & V_{2,4} \arrow[r] & V_{2,4}/aV_{2,3}
\end{tikzcd}
\]
vanishes. For fixed $a\ne0$ this is a nontrivial linear condition on $\delta$. In fact, choosing $\delta=x_0\partial_{x_0}+x_1\partial_{x_1}$,  one sees that it acts on $V_{2,4}$ as multiplication by $4$ and then does not
map all of $V_{2,4}$ into the proper subspace $aV_{2,3}$. In other words $T_4$ is a proper closed subset of the irreducible space $\operatorname{End}(V_{2,1})\times
\bigl(V_{2,1}\setminus\{0\})$. Hence its dimension is $\leq 5$. Thus $\operatorname{dim}(S_4) \leq 8$, and 
\begin{align*}
   \dim\{(g,h,D) \in \mathcal{I}_{3,4} \mid D\in S_4\}
   &\le\dim V_{3,4}+8-4\\
   &=\dim V_{3,4}+4=19.
\end{align*}
Now suppose \(a=0\) and \(\delta\neq0\). We claim that
\[
\rank\bigl(\delta\colon V_{2,4}\longrightarrow V_{2,4}\bigr)\ge4.
\]
As usual, recall that  $\delta  \in \operatorname{End}(V_{2,1}) \simeq \operatorname{Der}_{\CC}(V_2)_0$ (via the isomorphism \eqref{iso derivation}). Here we are considering the restriction of its action on $V_{2,4}$. To understand its  rank we relate the eigenvalues on  $V_{2,1}$ to the ones in $V_{2,4}$.  By the Jordan normal form, we may choose a basis \(u,v\) of
\(V_{2,1}\) such that
\[
\delta(u)=\alpha u,
\qquad
\delta(v)=\gamma u+\beta v.
\]
Here  \(\alpha\) and \(\beta\) are its eigenvalues, counted with
multiplicity. Using \eqref{iso derivation} one sees that  the associated derivation is determined by the Leibniz
rule
\[
D_\delta(\ell_1\ell_2\ell_3\ell_4)
=
\sum_{j=1}^4
\ell_1\cdots\ell_{j-1}\delta(\ell_j)
\ell_{j+1}\cdots\ell_4.
\]
Applying this formula to the monomial \(u^{4-i}v^i\), for
\(0\le i\le4\), gives
\[
\begin{aligned}
\delta|_{V_{2,4}}(u^{4-i}v^i)
&=(4-i)u^{3-i}v^i\delta(u)
  +i u^{4-i}v^{i-1}\delta(v)\\
&=(4-i)\alpha u^{4-i}v^i
  +i u^{4-i}v^{i-1}(\gamma u+\beta v)\\
&=\bigl((4-i)\alpha+i\beta\bigr)u^{4-i}v^i
  +i\gamma u^{5-i}v^{i-1}.
\end{aligned}
\]
Therefore, with respect to the ordered basis $\{u^4, u^3v, u^2v^2, uv^3,v^4\}$ of $V_{2,4}$, the matrix of $\delta|_{V_{2,4}}$ is upper triangular, and its eigenvalues are
\[
4\alpha,\qquad
3\alpha+\beta,\qquad
2\alpha+2\beta,\qquad
\alpha+3\beta,\qquad
4\beta.
\]
Thus
\[
   \det(\delta|_{V_{2,4}})
   =\prod_{i=0}^4\bigl((4-i)\alpha+i\beta\bigr).
\]
If $(\alpha,\beta)\ne(0,0)$, at most one factor in this product vanishes, so $\delta|_{V_{2,4}}$ has rank $4$ or $5$. It remains to consider the case $
\alpha=\beta=0.$ Since \(\delta\neq0\), the endomorphism \(\delta\) is nonzero
nilpotent. By the Jordan normal form, after rescaling \(v\), we may
assume
\[
\delta(u)=0,
\qquad
\delta(v)=u.
\]
The preceding formula then becomes
\[
\delta|_{V_{2,4}}(u^{4-i}v^i)
=
i\,u^{5-i}v^{i-1}.
\]
For \(i=0\), the image is zero, while the four images for $1\le i\le4$ are linearly independent. Hence
$\rank\delta|_{V_{2,4}}=4$.  We conclude that, for every nonzero
\(\delta\in\operatorname{End}(V_{2,1})\), $\rank\bigl(\delta|_{V_{2,4}})\ge4$ as desired. Therefore
\begin{align*}
   \dim\{(g,h,D) \in \mathcal{I}_{3,4}\mid a_D=0,\ \delta_D\ne0\}
   &\le\dim V_{3,4}+7-4\\
   &=\dim V_{3,4}+3=18.
\end{align*}
Finally, the locus $a=\delta=0$ has dimension $3$ and $r(0,0)=0$, so its corresponding incidence space also has dimension at most $\dim V_{3,4}+3=18$.
\end{proof}

\end{document}
